\clearpage
\section{Standalone calculation: complete setup and run guide}
\label{sec:run-guide}

This recipe starts a fresh quadrupole calculation at $T/T_c=0.30$,
$F_1^s=0$, without depending on an earlier run directory. It uses the full 2D
solver behind the historically named axisymmetric benchmark driver. All
relative paths below are resolved from the \emph{repository root}.

\subsection{Software and input inventory}

Required software is a Fortran compiler (the tested path is GNU Fortran),
CMake, a matching MPI library/launcher and Python with NumPy, Matplotlib and
Pillow. No OpenBLAS link path is needed for this workflow. On another machine,
build afresh rather than copying the laptop executable or its CMake cache.
For a Python virtual environment, for example:
\begin{lstlisting}
python3 -m venv .venv
source .venv/bin/activate
python -m pip install -r requirements.txt Pillow
\end{lstlisting}

The complete supplied namelist is
\code{examples/2d_quadrupole_Fs1_0.nml}. Its four auxiliary input files are
committed under
\begin{lstlisting}
benchmarks/a_phase_core_2d/reference/current_new_src_T0.30_F1s5.4/
    op_xyz
    curr
    gauss11.dat
    ozaki.dat
\end{lstlisting}
The driver currently requires the radial pair \code{op_xyz}/\code{curr}
even for a fresh non-radial seed. These files supply the comparison profile,
\emph{not} the quadrupole initial state or its fitted exterior. Their
$F_1^s=5.4$ label does not override \code{feedback_scale=0}. Radial comparison
errors and \code{benchmark passed} are not acceptance criteria for a
quadrupole at different parameters. Removing this compatibility dependency
from a generic production driver remains future work.

{\small
\begin{longtable}{p{31mm}p{109mm}}
\toprule
Input & Contents and authority\\
\midrule
\endhead
Namelist & Chooses mesh, active mask, seed, feedback, endpoint policy,
transport steps, iteration tolerance, checkpoints and file paths.\\
\code{op_xyz} & One radial coordinate followed by 18 real values for the
nine complex Cartesian order-parameter components, in row-major component
order $xx,xy,xz,yx,yy,yz,zx,zy,zz$, real/imaginary pairs.\\
\code{curr} & Five columns: radial coordinate, density diagnostic, and
three real current-related mean-field values. The radial importer uses the
fourth column (azimuthal field on the positive $x$ axis) and requires the
radial and axial fields to be negligible. Both radial files must describe
the same grid and cover the active comparison radius.\\
\code{gauss11.dat} & Eleven rows of polar node and Gauss weight. This driver
currently fixes 48 azimuthal samples and 11 polar nodes: 528 directions.
There is no namelist angular-count override in this driver.\\
\code{ozaki.dat} & Header: pole count, reduced temperature, cutoff;
then pole/residue pairs. The supplied file has 8 poles and $T/T_c=0.30$.
It sets the transport temperature and the seed temperature for this driver.
Changing temperature requires a consistent new table, not only a label or
the header of an old table.\\
Restart file & Required only for \code{restart}/\code{restart_core}.
For a fresh historical seed it is an empty string.\\
\bottomrule
\end{longtable}
}

\subsection{Complete fresh-run namelist}

This listing is imported directly from the runnable example when the note is
built, so the documented values do not drift from the supplied input.
\lstinputlisting{../../examples/2d_quadrupole_Fs1_0.nml}

\clearpage
\subsection{Meaning of the controls}

\begin{longtable}{>{\raggedright\arraybackslash\small}p{58mm}p{82mm}}
\toprule
Control or group & Interpretation and restrictions\\
\midrule
\endhead
\code{half_width}, \code{number_of_cells}, \code{smooth_stretch} &
$H=22$, $N=64$, $\lambda=2.3$ in \eqref{eq:smoothmesh}: 65 nodes per
axis, 4225 stored points, minimum spacing about $0.321\xi_0$. Use even $N$
to include the origin. These nodes are not moving during relaxation.\\
\code{active_radius} & Radius of independently updated points: 20 here,
3481 targets. Halo nodes outside are regenerated, not independently mixed.\\
\code{feedback_scale} & $F_1^s/(1+F_1^s/3)$, not $F_1^s$ itself.
Zero removes feedback in the trajectory self-energy; current-related
diagnostic values can still be computed and nonzero.\\
\code{asymptotic_bulk_gap}, \code{asymptotic_winding} &
Bulk amplitude 0.2799 in the current normalization, unit phase winding.
Changing temperature also requires a consistent bulk gap. Historical
presets assume a centred, singly quantized vortex; leave the centre at zero
and winding at one for this recipe.\\
\code{endpoint_policy} & Use \code{state_asymptotic} for tails determined
by the present relaxed state. \code{radial_reference} and
\code{radial_asymptotic} instead use prescribed radial information;
they are not substitutes for a freely evolving quadrupole exterior.\\
\code{asymptotic_fit_inner_radius}, \code{asymptotic_matching_radius} &
14 and 17 here. The order parameter uses the matching surface; the
mean field uses both. Ensure all bilinear fit stencil nodes are active,
not merely that the circles lie within the active radius.\\
\code{asymptotic_outer_radius} & Riccati endpoint circle, 70 here; it must
enclose the storage rectangle (for a centred square, $R_{\rm out}>\sqrt2H$).
It is not the computational cell radius.\\
\code{trajectory_maximum_step}, \code{asymptotic_maximum_step} &
Maximum spacing in path coordinate $s$ in the stored mesh and analytic
extension, respectively. Both 0.25 here. The runner's
\code{--trajectory-step} overrides only the first.\\
\code{minimum_substep_count}, \code{boundary_relaxation_distance} &
Floor for pole-dependent RK substeps and length of a constant endpoint
relaxation segment. Retained at 1 and 1.75.\\
\code{initialization_mode} & Physical choice of seed, embedding or restart.
Select one of the four historical presets in Section~\ref{sec:seeds}.\\
\code{maximum_iterations}, \code{convergence_tolerance} &
50 and $2\times10^{-6}$. Zero iterations requests a single map, not a
relaxation. The stopping tolerance applies to the largest absolute packed
real-component residual, not the area-weighted relative norm.\\
\code{anderson_history_limit}, \code{anderson_progress_threshold},
\code{anderson_maximum_mixing} & History capacity 10, source progress control
0.1 and mixing cap 5. Reaching the cap or filling history does not disable
the least-squares search.\\
\code{checkpoint_interval} & Save every accepted update here (1).
The field is saved but Anderson history is not.\\
\code{require_convergence}, \code{maximum_embedding_error} &
Benchmark acceptance controls, not modifications of the physical map.
False permits iteration-limited exploratory runs; the embedding threshold
checks the reference conversion, not physical agreement of a different core.\\
\bottomrule
\end{longtable}

For \code{extended_smooth}, additionally set \code{preserved_core_half_width}
and \code{outer_spacing_growth}; \code{half_width} becomes the new outer
half-width while \code{number_of_cells} describes the preserved base grid.
The optional \code{restart_core_radius} is used only by core transfer.
For \code{multiscale}, the fine/medium extents and three spacings determine
the grid instead. Fine/medium region controls also define diagnostic zones;
if omitted, the circular driver uses radii 2 and 5 for those zones. They do
not change a \code{smooth} grid's coordinates.

\subsection{Build and launch}

The Python runner configures a Release CMake build, builds the executable,
creates a uniquely dated run directory, writes the effective namelist,
launches MPI and plots saved output. From the repository root:
\begin{lstlisting}
python3 tools/run_axisymmetric_core_benchmark.py \
  --input examples/2d_quadrupole_Fs1_0.nml \
  --initialization-mode historical_qop \
  --initialization historical-qop \
  --case-name quadrupole-T030-Fs0 \
  --build-dir work/quadrupole-seed-build \
  --ranks 10 --formats png
\end{lstlisting}
\textbf{Important:} \code{--initialization} is a naming label only.
\code{--initialization-mode} or the namelist actually selects the seed.
The command specifies both to avoid ambiguity. To start a different core,
change the physical mode to \code{historical_nop}, \code{historical_aop}
or \code{historical_dop} and give the run a matching name. Do not assume
the compact quadrupole domain is validated for a swollen double core.

On the current Mac, the tested single-node launcher workaround prefixes that
command with
\begin{lstlisting}
caffeinate -i env OMPI_MCA_pml=ob1 OMPI_MCA_btl=self,sm
\end{lstlisting}
Put the prefix and Python command on the same shell command line.
This Open MPI 5 shared-memory setting is for this laptop, \emph{not} a
multi-node cluster configuration. On clusters use the site's compiler/MPI
modules and scheduler launcher, and omit \code{caffeinate} and this transport
override. Do not run two ten-rank jobs at the same time on the laptop.

For a manually managed build or scheduler job, the equivalent build target is
\begin{lstlisting}
cmake -S . -B work/standalone-build \
  -DCMAKE_Fortran_COMPILER=gfortran -DCMAKE_BUILD_TYPE=Release
cmake --build work/standalone-build \
  --target benchmark_axisymmetric_core_2d --parallel 10
\end{lstlisting}
The raw executable takes seven positional arguments: namelist, initial
field, mapped field, metrics, final field, history and checkpoint. Create
the output directory first. For example, use a unique directory rather than
overwriting a previous run:
\Needspace{10\baselineskip}
\begin{lstlisting}
run_dir="$(mktemp -d "$PWD/runs/manual-quadrupole-XXXXXX")"
cp examples/2d_quadrupole_Fs1_0.nml "$run_dir/input.nml"
mpiexec -n 10 work/standalone-build/benchmark_axisymmetric_core_2d \
  "$run_dir/input.nml" "$run_dir/input_fields_2d.dat" \
  "$run_dir/mapped_fields_2d.dat" "$run_dir/metrics.txt" \
  "$run_dir/final_fields_2d.dat" "$run_dir/iteration_history.dat" \
  "$run_dir/checkpoint_fields_2d.dat" > "$run_dir/run.log" 2>&1
\end{lstlisting}
This raw route does not automatically generate plots or the Python runner's
manifest. It still uses repository-root-relative input paths. For portable
archiving, copy all four auxiliary files into the archived case and update
their namelist paths explicitly; the runner archives the effective namelist
but does not make every referenced external file self-contained.

\subsection{Inspect, continue and interpret a run}

After the first checkpoint is written, a second terminal can produce an
independent inspection snapshot without evaluating any trajectories:
\begin{lstlisting}
python3 tools/plot_running_axisymmetric.py \
  --latest quadrupole --no-radial-comparison
\end{lstlisting}
For an older run, replace \code{--latest quadrupole} by its explicit directory.
Each invocation creates a dated inspection subdirectory: Cartesian and
harmonic amplitude/phase maps, pair-density and current-related proxies,
axis profiles, convergence and spatial RMS/max residual dynamics. It neither
stops the solver nor overwrites its files. If a checkpoint is being written,
retry after the write completes. The no-comparison flag avoids treating the
unrelated A-phase reference as the quadrupole target.

To continue from a checkpoint, substitute the actual directory name for the
dated placeholder below:
\begin{lstlisting}
old=runs/YYMMDD-HHMMSS-quadrupole-T030-Fs0-historical-qop
python3 tools/run_axisymmetric_core_benchmark.py \
  --input "$old/input.nml" \
  --restart "$old/checkpoint_fields_2d.dat" \
  --case-name quadrupole-continued --initialization restart \
  --build-dir work/quadrupole-seed-build \
  --ranks 10 --max-iterations 50 --checkpoint-interval 1 --formats png
\end{lstlisting}
Use \code{final_fields_2d.dat} after a completed run. Do not pass
\code{--initialization-mode historical_qop} when restarting: it would
override the requested restart. The mesh must match; a restart begins fresh
Anderson history and its iteration numbering restarts at one.

The live log prints one line per completed update, not continuous progress
within a map. A saved checkpoint after update $k$ contains the accepted field;
residual snapshot $k$ describes the field \emph{before} that update.
\code{iteration_limit} is not convergence, and the final accepted field need
not have its own newly evaluated residual. The plotting snapshot's internal
filename \code{final_fields_2d.dat} is only a compatibility copy of the
checkpoint, not a declaration of convergence.

Per-rank files ending in \code{shutdown.rankNNNNNN.txt} audit MPI finalization.
The observed launcher improper-exit warning is not proved to be a numerical
failure or harmless: retain the audit and exit status. Saved-output plotting
can proceed independently, but should never turn a failed launcher status
into a successful run. The current runner attempts plotting after launcher
failure when complete outputs exist and preserves that failure in its report.

\subsection{Current limits of this standalone route}

This is a runnable free-vortex calculation, not yet a general material/device
input language. There is no boundary-wall or lead specification in this
namelist, no connected dipole self-energy, no arbitrary exchange-field input,
and no GPU backend. The driver exposes circular masks and theoretical bulk
orientation; optional zero-mode projection, elliptical mask studies and
shadow-probe workflows are in the separate double-core driver, not controls
that can be copied into this namelist. \code{prepare_only=.true.} is an
internal executable preflight that writes the initial/reference fields and
exits without maps; it does not produce the complete files expected by the
normal Python run wrapper. Keep it false or omitted for a calculation.
